% Related work. Citations: refs.bib (every entry checked by verify_refs.py).
% Source notes for each statement: related_work_notes.md.
\section{Related work}
\label{sec:related}

\subsection{Oxidation of refractory complex concentrated alloys}

Refractory high-entropy and complex concentrated alloys (RHEAs/RCCAs) keep much of their strength
to very high temperatures \citep{senkov2018review,senkov2019high}, but most of them oxidise as
badly as the refractory metals they contain: they pest, spall, lose mass by forming volatile oxides,
or oxidise completely \citep{gorr2021review}. The ductile Hf--Nb--Ta--Ti--Zr alloys, for example,
pest between 600 and 1000\,\textdegree C \citep{sheikh2018pesting}. Protective behaviour is found in
narrow composition windows. Alloys of the Al--Cr--Mo--Ta--Ti family form a slow-growing CrTaO$_4$
layer \citep{gorr2020crtao4,lo2019crtao4}. Replacing Ta by Nb gives porous, spalling scales
\citep{muller2019mechanism}, and the Al content decides whether a dense CrTaO$_4$ layer forms or
volatilisation of MoO$_3$ and CrO$_3$ dominates \citep{schellert2021al}. Two consequences matter
for data-driven models. First, the property changes sharply with composition, so the accuracy of a
model should depend on how close a query alloy is to measured alloys. Second, the usual target,
net specific mass change, combines scale growth with volatilisation and spallation, and its kinetics
change with exposure time \citep{lo2019crtao4}. Values from different sources are therefore noisy
and not always comparable.

\subsection{Public data}

Two public compilations of RCCA oxidation data were built independently. The dataset of
\citet{lin2026dataset} (886 records from 62 publications and in-house tests; the data used by
\citet{gorsse2025ml}) was extracted independently by two curators, conflicts were traced back to the
sources, and suspected factor-of-ten errors from unit conversion or transcription were checked and
removed when they could not be confirmed. RefOxDB \citep{mishra2024refoxdb} holds 407 mass-change
curves digitised from 68 sources. Its authors note that units of mass gain vary between papers, and
they find that the kinetic law chosen by the original authors differs from the best-supported law
in 71\% of the curves. Both groups digitised plotted data with WebPlotDigitizer
\citep{rohatgi2018webplotdigitizer}. Studies in other fields show that two people digitising the same
figure agree closely \citep{drevon2017intercoder}, that reliability depends on the figure rather than
on the program \citep{moeyaert2016reliability}, and that close agreement between people can coexist
with errors of up to 30\% against the true values \citep{burda2017estimating}. Automated extraction
of materials data from text reaches a precision of 73\% \citep{court2018auto} to about 90\%
\citep{polak2024extracting}. We did not find a study that compares two independent digitisations of
the same materials-property figures. Because the two oxidation compilations overlap in source papers,
they make this comparison possible, and they also make a true external test possible: a model fitted
to one compilation can be tested on the alloys that only the other contains.

High-temperature strength data come from separate compilations
\citep{couzinie2018comprehensive,gorsse2018database,borg2020expanded}, which contain no oxidation
data. \citet{borg2020expanded} also corrected errors in the earlier compilations, a reminder that
curated datasets need auditing.

\subsection{Data-driven oxidation models for RCCAs}

\citet{gorsse2025ml} trained gradient-boosted trees on the dataset of \citet{lin2026dataset} to predict
specific mass gain from composition, temperature and time. They validated the model by nested k-fold
cross-validation with shuffling over records, where a record is one combination of alloy,
temperature and time, so records of the same alloy can fall in both the training and the test folds.
They also tested five new alloys (25 records), for which the error was slightly higher than in
cross-validation, and three designed Al--Cr--Mo--Ta--Ti alloys, whose measured mass gains were close
to the targets. \citet{bejjipurapu2026descriptors} fitted a Gaussian process with oxidation-product
descriptors to 77 compositions from RefOxDB (24\,h at 900--1000\,\textdegree C). The test error on a
random 20\% split was 5.78\,mg/cm$^2$, while repeated five-fold cross-validation gave about
11\,mg/cm$^2$, which shows how much a single split can vary at this data size. The same group then ran
an experimental active-learning campaign from 81 literature alloys and found two alumina-forming
Al--Cr--Mo--Ti alloys with less than 1\,mg/cm$^2$ after 24\,h at 1000\,\textdegree C
\citep{bejjipurapu2025active}. Their reported predictions also show how far the model could be wrong
away from its data: several alloys predicted at 1--3\,mg/cm$^2$ gained 70--121\,mg/cm$^2$, and the two
five-component alloys with the lowest predictions (0.79 and 0.93\,mg/cm$^2$) gained 4.8 and
5.9\,mg/cm$^2$ \citep[SI Table 2]{bejjipurapu2025active}. Those alloys were chosen by expected
improvement, which deliberately samples uncertain regions, so these misses are consistent with, but do
not by themselves demonstrate, the selection bias discussed below. Data-driven oxidation models also
exist for Ti alloys \citep{bhattacharya2020ti}, Ni-based superalloys \citep{kim2021ni}, NiCr alloys
\citep{peng2021nicr}, several alloy classes together \citep{taylor2021ml} and Ni--Co--Cr--Al--Fe
high-entropy alloys \citep{tan2025nicocralfe}. \citet{peng2021nicr} attribute their success to a
consistent dataset from one laboratory, which the RCCA compilations are not.

To our knowledge, no RCCA oxidation study groups validation folds by source publication or tests a
model fitted to one compilation on the other.

\subsection{Strength models and multi-property design}

The edge-dislocation theory of \citet{maresca2020mechanistic} predicts the yield strength of BCC
high-entropy alloys without fitted parameters. It reproduces the strength of MoNbTaW and MoNbTaVW
up to 1900\,K, and its reduced form was used to screen the Mo--Nb--Ta--V--W family.
\citet{lee2021strength} showed experimentally that edge dislocations also control the strength of
NbTaTiV and CrMoNbV, and \citet{baruffi2022screw} found that the edge theory agrees with experiment
only in non-dilute alloys with a sufficiently large misfit parameter; screw dislocations control
strength otherwise. The theory is therefore supported for non-dilute, high-misfit BCC alloys, and
using it for the Al-containing compositions that resist oxidation is an extrapolation.
\citet{senkov2019high} related the loss of strength of single-phase BCC RCCAs to diffusion-controlled
deformation, which starts at about 0.6\,$T_\mathrm{m}$, and found that multi-phase RCCAs lose strength
rapidly above 1000\,\textdegree C; specific strength at 1000\,\textdegree C is therefore a demanding
objective.

Multi-property design of refractory alloys has combined the edge theory with phase, density, solidus
and ductility constraints \citep{vela2023high}, fused the edge theory with the data of
\citet{borg2020expanded} into a Bayesian strength model \citep{vela2023data}, optimised ductility
indicators with Bayesian optimisation under learned constraints \citep{khatamsaz2022multi}, and added
uncertainty quantification to strength--plasticity trade-offs \citep{giles2026yield}. Phase
constraints in such searches usually come from classifiers whose cross-validated accuracy is
moderate: 64--74\% for three phase classes \citep{huang2019machine} and 93\% for solid-solution
formation \citep{pei2020machine}. None of these studies includes oxidation. \citet{gorsse2025ml}
propose combining their oxidation model with strength and ductility models in a Pareto analysis,
which is the use case examined here.

\subsection{Validation beyond random splits}

When data are clustered or redundant, random cross-validation overstates the error on new cases. This
has been shown for ecological data \citep{roberts2017cross}, for compound activity prediction, where
time splits match prospective performance better than random splits \citep{sheridan2013time}, and
across machine-learning-based science generally, where leakage between training and test data
affected hundreds of papers \citep{kapoor2023leakage}. In materials science, \citet{meredig2018can}
introduced leave-one-cluster-out cross-validation, \citet{xiong2020evaluating} proposed forward
cross-validation to measure extrapolation to extreme property values, and \citet{omee2024structure}
showed that graph neural networks lose much of their accuracy under out-of-distribution splits.
Large materials datasets are so redundant that most of their data can be removed with little effect on
in-distribution accuracy, but not on out-of-distribution accuracy \citep{li2023exploiting}.
\citet{afflerbach2022machine} grouped a multi-source experimental dataset by alloy system to validate
models of the critical cooling rate of metallic glasses.

Error grows with dissimilarity to the training data. This was first shown for compound activity
models \citep{sheridan2004similarity}, and has since been used as an uncertainty measure
\citep{janet2019quantitative} and to define domains of applicability \citep{sutton2020identifying,
schultz2025general}; \citet{schultz2025general} found that high dissimilarity goes with large errors
and unreliable uncertainty estimates. \citet{li2025probing} showed that many nominally
out-of-distribution tests in materials science are in fact interpolation, which argues for measuring
the distance to the training data directly instead of relying on the split label. Tests on new data
are rarer. Models trained on the 2018 Materials Project degraded severely on compounds added by 2021
\citep{li2023critical}, and medical imaging models validated internally performed worse at hospitals
not seen in training \citep{zech2018variable}. Good regression metrics also do not guarantee good
design decisions \citep{bartel2020critical,borg2023quantifying}. By contrast, models trained on
computed energies of simple alloys transfer well to more complex high-entropy compositions
\citep{li2024efficient}, which suggests that the data type matters. Most of this evidence comes from
computed data or from other fields. Experimental data compiled from many papers add effects of the
source, such as laboratory practice, units and digitisation, that a distance in composition space
cannot detect.

\subsection{Uncertainty and design under model error}

A predictive interval is calibrated if it contains the true value as often as its level states
\citep{kuleshov2018accurate}. Materials studies have proposed metrics for comparing uncertainty
estimates \citep{tran2020methods} and shown that rescaling the spread of an ensemble on held-out data
improves calibration \citep{palmer2022calibration}; for molecular properties no uncertainty method was
best across datasets \citep{hirschfeld2020uncertainty}. These studies check calibration on held-out
splits of the same dataset. The calibration factor used here is fitted in the same way, on grouped
folds, and is then tested on the independent compilation.

Selecting the best of many candidates by their estimated value makes the selected candidate's estimate
too optimistic on average, even when each estimate is unbiased: the optimiser's curse
\citep{smith2006optimizer}. In model-based design the same effect appears as an optimiser finding
inputs where the model is wrong in its favour, and proposed remedies keep the search near the training
data \citep{brookes2019conditioning}, retrain the model as the search moves \citep{fannjiang2020autofocused},
or fit models that are conservative away from the data \citep{trabucco2021conservative}. Trust regions
in Bayesian optimisation \citep{eriksson2019scalable} serve a different purpose: they make the search
efficient while new evaluations arrive. The trust region used here is a fixed constraint on distance to
the data and on predicted uncertainty, applied in an offline setting with no new measurements. In
alloy discovery, \citet{ren2018accelerated} found that an initial model predicted metallic-glass
formation in a new ternary system in good qualitative agreement with experiment but with quantitative
errors in composition, and that retraining on these discrepancies improved its accuracy.

Confidence intervals in this work resample whole alloys, source papers or compositions, because records
within these groups are not independent; the cluster bootstrap is consistent for such data
\citep{field2007bootstrapping}, but inference can be unreliable when the number of clusters is small
\citep{cameron2008bootstrap}.

\subsection{Contribution relative to prior work}

Compared with the work above, this paper
(i) audits the two public RCCA oxidation compilations against each other, including a $\times$100 unit
conflict that is visible only when the two are compared;
(ii) measures how accuracy depends on composition distance to the training alloys, both in grouped
cross-validation and on an independently compiled dataset;
(iii) shows, with bootstrap confidence intervals, how the estimated accuracy and the model ranking change
from record-level to alloy-level, source-level and external validation;
(iv) shows that strength and oxidation data share very few alloys, so joint design extrapolates; and
(v) measures how strongly the optimum found with one surrogate depends on that surrogate.
The general principles are known; the contribution is to quantify them for an experimental,
literature-compiled alloy property where design decisions depend on them.
